% ID: SEL-D0034-OB-E02
% ESTADO: validada. S035, cotejada con examen y criterios.
\textbf{a)} En $2$ el límite izquierdo es $a-4$ y el valor
y límite derecho son $1$. La continuidad exige $\boxed{a=5}$.
Las derivadas de los tramos son $2x-4$ y $-1/(x-1)^2$.
En el empalme dan pendientes laterales $0$ y $-1$, así que
no existe derivada en $2$.

\textbf{b)} Para $a=4$ el primer tramo es $(x-2)^2$, con
derivada negativa en $(-\infty,2)$. El segundo decrece en
$[2,\infty)$. Al pasar por $2$, el límite izquierdo es $0$
y $f(2)=1$, por lo que hay un salto ascendente. La función
decrece en cada tramo, pero no en todo su dominio.

A la derecha, $1/(x-1)\to0$, de modo que $y=0$ es asíntota
horizontal. El denominador no se anula en $x\ge2$ y el salto
tiene límites finitos, así que no hay asíntotas verticales.
Por la izquierda el crecimiento cuadrático descarta cualquier
asíntota recta.
